% Part: many-valued-logic
% Chapter: three-valued-logics
% Section: goedel

\documentclass[../../../include/open-logic-section]{subfiles}

\begin{document}

\olfileid[tr]{mvl}{thr}{god}

\olsection{G\"odel Mantıkları}

Kurt G\"odel, her biri sezgisel mantıkta geçerli bütün !!{formula}s{}i
içeren ve klasik mantığın içinde yer alan bir $n$-değerli mantıklar dizisi
ortaya koydu. Bunların ilk ilginç örneği şöyledir:

\begin{defn}\ollabel{defn:goedel}
\emph{$3$-değerli G\"odel mantığı}~$\LogGod$ şu matrisle tanımlanır:
\begin{enumerate}
  \item $\lfalse$, $\lnot$, $\land$, $\lor$, $\lif$ bağlaçlarını içeren
  standart önerme dili $\Lang L_0$.
  \item Doğruluk değerleri kümesi $V = \{\True, \Undef, \False\}$.
  \item Tek belirlenmiş değer $\True$'dur; yani $V^+ = \{\True\}$.
  \item $\lfalse$ için $\tf{\lfalse} = \False$'dur. Geri kalan bağlaçların
  doğruluk fonksiyonları aşağıdaki tablolarla verilir:
  \begin{center}
    \begin{tabular}{c|c} 
      $\tf{\lnot}[\LogGod]$ & \\ 
      \hline  
      $\True$ & $\False$ \\ 
      $\Undef$ & $\False$ \\
      $\False$ & $\True$ 
    \end{tabular}
    \quad
    \begin{tabular}{c|ccc} 
      $\tf{\land}[\LogGod]$ & $\True$ & $\Undef$ & $\False$ \\ 
      \hline 
      $\True$ & $\True$ & $\Undef$ & $\False$ \\ 
      $\Undef$ & $\Undef$ & $\Undef$ & $\False$\\ 
      $\False$ & $\False$ & $\False$ & $\False$ 
    \end{tabular}
    \\[2ex]
    \begin{tabular}{c|ccc} 
      $\tf{\lor}[\LogGod]$ & $\True$ & $\Undef$ & $\False$ \\ 
      \hline 
      $\True$ & $\True$ & $\True$ & $\True$ \\ 
      $\Undef$ & $\True$ & $\Undef$ & $\Undef$ \\
      $\False$ & $\True$ & $\Undef$ & $\False$ 
    \end{tabular}
    \quad
    \begin{tabular}{c|ccc} 
      $\tf{\lif}[\LogGod]$ & $\True$ & $\Undef$ & $\False$ \\ 
      \hline 
      $\True$ & $\True$ & $\Undef$ & $\False$ \\ 
      $\Undef$ & $\True$ & $\True$ & $\False$  \\ 
      $\False$ & $\True$ & $\True$ & $\True$ 
    \end{tabular}
  \end{center} 
\end{enumerate}
\end{defn}

$\land$ ve~$\lor$ doğruluk tablolarının \L ukasiewicz ve güçlü Kleene
mantığındakilerle aynı, $\lnot$ ve~$\lif$ tablolarının ise üçünde de farklı
olduğuna dikkat edin. G\"odel mantığında
$\tf{\lnot}(\Undef) = \False$'dur. \L ukasiewicz ve Kleene mantığından farklı
olarak $\tf{\lif}(\Undef, \False) = \False$; Kleene mantığından farklı
(fakat \L ukasiewicz mantığındaki gibi) olarak
$\tf{\lif}(\Undef, \Undef) = \True$'dur.

Yukarıda değinilen sezgisel mantık bağlantısının düşündürdüğü gibi
$\LogGod[3]$, sezgisel mantığa yakındır. Sezgisel mantığın bütün doğruları
$\LogGod[3]$ içinde totolojidir; sezgisel olarak geçerli olmayan birçok
klasik totoloji de $\LogGod[3]$ içinde totoloji değildir. Örneğin
aşağıdakiler totoloji değildir:
\begin{align*}
  & p \lor \lnot p && (p \lif q) \lif (\lnot p \lor q) \\
  & \lnot\lnot p \lif p && \lnot(\lnot p \land \lnot q) \lif (p \lor q) \\
  & ((p \lif q) \lif p) \lif p && \lnot(p \lif q) \lif (p \land \lnot q)
\end{align*}
Ancak $\LogGod[3]$'ün her totolojisi sezgisel olarak geçerli değildir;
örneğin $\lnot\lnot p \lor \lnot p$ ya da
$(p \lif q) \lor (q \lif p)$.

\begin{prob}
  Aşağıdakilerin $\LogGod[3]$'ün totolojileri olduğunu gösterecek doğruluk
  tabloları verin:
  \begin{align*}
    & \lnot\lnot p \lor \lnot p\\
    & (p \lif q) \lor (q \lif p) \\
    & \lnot(p \land q) \lif (\lnot p \lor \lnot q) \\
    & (p \lif q) \lor (q \lif r) \lor (r \lif s)
  \end{align*}
\end{prob}

\begin{prob}
  Aşağıdakilerin $\LogGod[3]$'ün totolojileri olmadığını gösterecek doğruluk
  tabloları verin:
  \begin{align*}
    & (p \lif q) \lif (\lnot p \lor q) \\
    & \lnot(\lnot p \land \lnot q) \lif (p \lor q) \\
    & ((p \lif q) \lif p) \lif p \\
    & \lnot(p \lif q) \lif (p \land \lnot q)
  \end{align*}
\end{prob}

\begin{prob}
  Aşağıdaki bağıntılardan hangileri G\"odel mantığında geçerlidir? Her biri
  için bir doğruluk tablosu verin.
  \begin{enumerate}
    \item $p, p \lif q \Entails q$
    \item $p \lor q, \lnot p \Entails q$
    \item $p \land q \Entails p$
    \item $p \Entails p \land p$
    \item $p \Entails p \lor q$
  \end{enumerate}
\end{prob}

\end{document}
